\section{Modelos del mundo, tacto y manos diestras}

Las secciones anteriores trataron la transferencia de acciones y los datos; esta sección aborda dos capacidades más fundamentales: los modelos del mundo y el tacto. Tan Jie (谭杰) describe el modelo del mundo como Vision-Language-Vision: recibe la visión actual y el lenguaje o la acción, y genera el siguiente fotograma o el estado visual futuro. El tacto, por su parte, cobra importancia en el escenario de las manos diestras, porque el agarre fino, el deslizamiento, el material y la fuerza de contacto no pueden resolverse de forma estable solo con la visión.

\begin{figure}[H]
\centering
\includegraphics[width=0.96\textwidth]{world-model-vlv.png}
\caption{Modelo del mundo V-L-V: entrada de Vision + Language, predicción del estado visual del siguiente fotograma. Diagrama conceptual propio, elaborado a partir de la conversación en 01:03:48--01:08:29.}
\end{figure}

\begin{knowledgebox}{Lectura de la figura: el modelo del mundo «predice las consecuencias de las acciones»}
La imagen actual y el objetivo expresado en lenguaje entran al World Model; el modelo predice el siguiente fotograma o el estado visual futuro, y con base en ello la Policy elige la acción. No se limita a generar videos atractivos, sino que ayuda al robot a imaginar «si hago esto, ¿qué pasará con el mundo?».
\end{knowledgebox}

\subsection{Por qué importan los modelos del mundo}

Esta subsección lleva el modelo del mundo de la «generación de video» de vuelta al control de robots. Lo que el robot necesita es predecir las consecuencias de sus acciones, no generar una secuencia de imágenes que parezca verosímil.

El robot necesita predecir las consecuencias antes de actuar. Si al tomar una taza esta se volcará, en qué dirección hay que jalar un cierre, cómo se deformará una tela al tocarla con la mano: todo esto requiere modelar los cambios en el estado del mundo. Un modelo del mundo suficientemente capaz puede reducir el costo del ensayo y error en el mundo real y mejorar la capacidad de planificación y de recuperación.

\subsection{Por qué el tacto vuelve a ser importante}

A continuación se analizan los cambios en el nivel de los sensores. Antes se solía considerar que el tacto no era importante, en parte porque el hardware no había llegado a una etapa que exigiera un tacto fino; en cuanto se pasa a las manos diestras, el tacto deja de ser un complemento opcional y se convierte en una entrada básica.

Tan Jie menciona que, si antes se pensaba que el tacto no era importante, era por las limitaciones del hardware; con una mano diestra, el tacto se vuelve muy importante. Una pinza común puede completar muchas tareas con visión y movimientos gruesos; una mano diestra, en cambio, necesita conocer la fuerza de contacto, el deslizamiento, el material, la estabilidad del agarre y el estado de las yemas de los dedos. Sin tacto, mucha manipulation fina difícilmente se completa de forma confiable.

\begin{figure}[H]
\centering
\includegraphics[width=0.94\textwidth]{tactile-dexterous-hand.png}
\caption{Tacto y manos diestras: cuanto más diestro es el hardware, más importante es el tacto. Diagrama conceptual propio, elaborado a partir de la conversación en 01:08:29--01:17:35.}
\end{figure}

\begin{knowledgebox}{Lectura de la figura: el valor del tacto aumenta con la complejidad del cuerpo del robot}
Una pinza común se ocupa principalmente de la posición y del cierre; una mano diestra debe manejar la coordinación de varios dedos, la fuerza de contacto, el deslizamiento, el material y la estabilidad del agarre. Cuanto más se parece el hardware a la mano humana, menos es el tacto un sensor prescindible.
\end{knowledgebox}

\subsection{Resumen de la sección}

El modelo del mundo resuelve la «predicción de consecuencias» y el tacto resuelve el «contacto fino». Ambos apuntan a lo mismo: para pasar de movimientos burdos a una manipulación confiable, el robot debe comprender los cambios y los contactos del mundo físico.
